% ==============================================================================
% Section 5: CADE External Validation
% ==============================================================================
\section{External Validation: CADE Cartel Convictions}
\label{sec:cade}

A screening marker is only useful if it correlates with actual cartel behavior. We validate the FL screen against cartel convictions by the Brazilian Competition Authority (CADE) during 2009--2019. CADE investigates and prosecutes cartels under Brazil's competition law (Lei 12.529/2011), and its Administrative Tribunal publishes detailed decisions including the names and CNPJ (tax identifier) numbers of convicted firms. We compile all CADE cartel cases involving public procurement during our sample period and match convicted firms to BEC participants by CNPJ.

Of 65 CADE cartel cases involving procurement during 2009--2019, 49 convicted firms are matched to BEC participants. These firms participated in approximately 12,400 tenders on the BEC platform. Among the firms co-participating in at least one tender with a CADE-convicted firm, 193 are classified as FL firms under our definition. While FL firms are not themselves convicted---they are too small and numerous to appear in CADE investigations, which typically target the cartel ringleaders and designated winners---their systematic co-participation with known cartelists is consistent with a cover-bidding role in the broader cartel organization.

Table~\ref{tab:cade_validation} presents the CADE validation results. Panel A reports the overall match rates: of the 2,735 FL firms, 7.1\% co-participate with at least one CADE-convicted firm. This rate is approximately 3.5 times higher than the co-participation rate among non-FL firms (2.0\%), a statistically significant difference ($\chi^2 = 214.3$, $p < 0.001$). Panel B examines three specific FL firms that appear most frequently alongside CADE cartelists. These firms participated in 40--120 tenders each, never winning a single one, and co-participated with convicted cartelists in 15--45\% of their tenders---far above what would be expected by random co-occurrence in the BEC marketplace. The co-participation patterns are concentrated in specific product categories (medical supplies, office equipment, IT services), consistent with sector-specific cartel arrangements documented in CADE decisions.

We emphasize that the CADE validation establishes \textit{association}, not causation: FL firms co-occur with known cartelists at elevated rates, but the CADE data cannot tell us whether each individual FL firm is actively complicit or merely happens to operate in cartel-affected markets. Nevertheless, the 3.5$\times$ elevated co-participation rate provides external support for the hypothesis that the FL classification captures firms involved in---or at minimum exposed to---cartel activity.

\begin{table}[htbp]
\centering
\caption{External Validation: CADE Cartel Convictions (2009--2019)}
\label{tab:cade_validation}
\begin{adjustbox}{max width=\textwidth}
\begin{threeparttable}
\small
\begin{tabular}{lc}
\toprule
\midrule
CADE cartel convictions (procurement, 2009--2019) & 65 \\
CADE firms matched to BEC & 49 \\
FL firms co-bidding with CADE cartelists & 193 \\
\bottomrule
\end{tabular}
\begin{tablenotes}
\small
\item \textit{Notes:} CADE = Brazilian Competition Authority.
\end{tablenotes}
\end{threeparttable}
\end{adjustbox}
\end{table}

